\SourceLeaf{307}
\Needspace{8\baselineskip}\begingroup\small
\begin{center}
BEMERKUNGEN.
\end{center}

Die Notiz \emph{Zur Theorie der Kegelschnitte} bezieht sich auf die im Brief vom 12. Mai 1843 erwähnten Untersuchungen. Sie ist aus dem Nachlass in der folgenden Abteilung Verwendung complexer Grössen für die Geometrie auf S. 341 bis 344 dieses Bandes abgedruckt.

\textsc{Gauss}' Lösung der von \textsc{Möbius} gestellten geometrischen Aufgabe hat \textsc{Schumacher} in den Astronomischen Nachrichten Nr. 42, 1823 November, veröffentlicht. Sie ist in Bd. IV, S. 406---407, wieder abgedruckt.

Die Notiz [II] steht auf einem einzelnen Blatt. Sie stammt wahrscheinlich aus derselben Zeit wie die unter [I] abgedruckten Briefstellen. \textsc{Gauss} hat den baryzentrischen Kalkül, wie unter [3] angedeutet, insbesondere auf die Theorie des Fünfecks angewandt. Die betreffenden Entwicklungen, die mit dem Pentagramma mirificum zusammenhängen, findet man unter dem Titel \emph{Weitere Fragmente über das Pentagramma mirificum}, Nr. 11, S. 109 bis 111 dieses Bandes (vgl. auch die zugehörigen Bemerkungen \textsc{Frickes}, S. 117). Hinzuweisen ist auch auf die Notiz \emph{Beziehung der Raumverhältnisse auf ein gegebenes Tetraeder}, Bd. II, S. 307 bis 308. Sie stammt wahrscheinlich schon aus dem Juli 1831.

\begin{flushright}
\textsc{Stäckel}.
\end{flushright}
\par\endgroup
